% Irodov, user-supplied Arihant PDF, section 2.2, original problems 46--56,
% printed pages 67--69 (PDF pp.76--78). One-mole expressions unless stated.
\thermoenergyproblem{46}{67}
氩气在一次多方过程中，体积增大到原来的
$\alpha=4.0$ 倍，同时压强降至原来的
$1/\beta$，其中 $\beta=8.0$。
把氩气视为理想单原子气体，求此过程的摩尔热容。
\exsolution
设多方过程为 $pV^n=\text{常数}$，
则
\[
\frac{p_2}{p_1}=\left(\frac{V_1}{V_2}\right)^n
\Longrightarrow \frac1\beta=\alpha^{-n},\qquad
n=\frac{\ln\beta}{\ln\alpha}
=\frac{\ln8}{\ln4}=\frac32.
\]
由理想气体第一定律，
过程摩尔热容
\[
C=C_V+\frac{R}{1-n}.
\]
单原子气体 $C_V=3R/2$，
故
\[
\boxed{C=\frac{3R}{2}+\frac{R}{1-3/2}
=-\frac R2\simeq-4.16\,\mathrm{J/(mol\,K)}}.
\]
\excheck 负热容是\emph{此多方路径的热量变化率}
$\delta Q/\mathrm dT<0$，不是单原子气体
平衡定容热容变成负值。

\thermoenergyproblem{47}{68}
一摩尔氩气沿多方指数 $n=1.50$ 的过程膨胀，
温度变化 $\Delta T=-26\,\mathrm K$。
求：(a) 气体吸热 $Q$；(b) 气体做功 $A$。
\exsolution
单原子氩气 $C_V=3R/2$，
多方过程的摩尔热容
\[
C=C_V+\frac R{1-n}
=\frac{3R}{2}-2R=-\frac R2.
\]
因此
\[
\boxed{Q=C\Delta T
=\left(-\frac R2\right)(-26\,\mathrm K)
\simeq+108\,\mathrm J}.
\]
理想气体沿多方过程有
$p\,\mathrm dV=R\,\mathrm dT/(1-n)$，
故
\[
\boxed{A=\frac{R\Delta T}{1-n}
=\frac{8.314(-26)}{-0.50}\,\mathrm J
\simeq432\,\mathrm J}.
\]
\excheck 内能变化
$\Delta U=(3R/2)(-26)\simeq-324\,\mathrm J$，
故 $\Delta U+A\simeq108\,\mathrm J=Q$，
与第一定律一致。

\thermoenergyproblem{48}{68}
绝热指数为 $\gamma$ 的理想气体
沿 $p=\alpha V$ 的规律膨胀，
其中 $\alpha$ 为常数，
初体积 $V_0$，末体积 $\eta V_0$。
求：(a) 内能增量；(b) 气体做功；
(c) 该过程的摩尔热容。
\exsolution
设气体含 $\nu$ 摩尔，则
$pV=\nu RT$、$\nu C_V=\nu R/(\gamma-1)$。
沿 $p=\alpha V$ 有
$\nu RT=\alpha V^2$，所以
\[
\Delta U=\nu C_V\Delta T
=\frac{\alpha}{\gamma-1}(V_2^2-V_0^2).
\]
取 $V_2=\eta V_0$，
\[
\boxed{\Delta U=
\frac{\alpha V_0^2(\eta^2-1)}{\gamma-1}}.
\]
机械功
\[
\boxed{A=\int_{V_0}^{\eta V_0}\alpha V\,\mathrm dV
=\frac{\alpha V_0^2}{2}(\eta^2-1)}.
\]
微分状态方程有
$\nu R\,\mathrm dT=2\alpha V\,\mathrm dV$，
故每摩尔过程热容
\[
C=C_V+\frac{p\,\mathrm dV}{\nu\,\mathrm dT}
=\frac R{\gamma-1}+\frac R2
=\boxed{\frac{R(\gamma+1)}{2(\gamma-1)}}.
\]
\excheck 题目未指定摩尔数；
$\Delta U,A$ 是\emph{整团气体的}能量，
而 $C$ 是每摩尔的热容。

\thermoenergyproblem{49}{68}
绝热指数为 $\gamma$ 的理想气体在膨胀中
吸收的热量恰好等于其\emph{内能减少的量}。
求：(a) 过程摩尔热容；(b) $T,V$ 过程方程；
(c) 一摩尔气体从初温 $T_0$ 出发、
体积增大 $\eta$ 倍时对外做功。
\exsolution
题意为 $\delta Q=-\mathrm dU$，
由第一定律
$\delta Q=\mathrm dU+p\,\mathrm dV$
得到
$p\,\mathrm dV=-2\,\mathrm dU=-2\nu C_V\,\mathrm dT$。
于是
\[
\boxed{C=\frac{\delta Q}{\nu\,\mathrm dT}=-C_V
=-\frac R{\gamma-1}}.
\]
由 $pV=\nu RT$，
\[
\frac{\mathrm dT}{T}
=-\frac{R}{2C_V}\frac{\mathrm dV}{V}
=-\frac{\gamma-1}{2}\frac{\mathrm dV}{V}.
\]
积分：
\[
\boxed{TV^{(\gamma-1)/2}=\text{常数}}.
\]
若为一摩尔，末温
$T_2=T_0\eta^{-(\gamma-1)/2}$，
机械功 $A=-2\Delta U$，故
\[
\boxed{A=\frac{2RT_0}{\gamma-1}
\left[1-\eta^{-(\gamma-1)/2}\right]}.
\]
\excheck 膨胀导致 $T_2<T_0$，
因此 $\Delta U<0$，
$\delta Q=-\mathrm dU>0$，即仍吸热。

\thermoenergyproblem{50}{68}
一摩尔绝热指数为 $\gamma$ 的理想气体，
经历过程 $p=aT^\alpha$，$a,\alpha$ 为常数。
求：(a) 温度增加 $\Delta T$ 时气体做功；
(b) 过程摩尔热容与热容为负的条件。
\exsolution
由理想气体状态方程
$V=RT/p=(R/a)T^{1-\alpha}$，
故
\[
\mathrm dV=(R/a)(1-\alpha)T^{-\alpha}\mathrm dT,
\qquad p\,\mathrm dV=R(1-\alpha)\mathrm dT.
\]
因此
\[
\boxed{A=R(1-\alpha)\Delta T}.
\]
再用
$\delta Q=C_V\,\mathrm dT+p\,\mathrm dV$，
\[
\boxed{C=C_V+R(1-\alpha)
=R\left(\frac{\gamma}{\gamma-1}-\alpha\right)}.
\]
当
\[
\boxed{\alpha>\frac{\gamma}{\gamma-1}}
\]
时 $C<0$。
\excheck $\alpha=1$ 是等容路径，
$A=0$；$\alpha=0$ 是等压路径，
$C=C_p$，均可反查。

\thermoenergyproblem{51}{68}
某理想气体的内能和体积沿过程中满足
$U=aV^\alpha$，
其中 $a,\alpha$ 为常数，
气体绝热指数为 $\gamma$。
若内能增量为 $\Delta U$，求：(a) 气体做功
和传递给气体的热量；(b) 过程摩尔热容。
\exsolution
理想气体满足
$pV=\nu RT=(\gamma-1)U$，
所以
\[
p=\frac{(\gamma-1)U}{V}
=(\gamma-1)aV^{\alpha-1}.
\]
对于 $\alpha\ne0$，
\[
A=\int p\,\mathrm dV
=\frac{\gamma-1}{\alpha}
\,a(V_2^\alpha-V_1^\alpha)
=\boxed{\frac{\gamma-1}{\alpha}\Delta U}.
\]
由第一定律
\[
\boxed{Q=\Delta U+A
=\frac{\alpha+\gamma-1}{\alpha}\Delta U}.
\]
因 $\mathrm dU=\nu C_V\,\mathrm dT$，
过程摩尔热容为
\[
\boxed{C=C_V\frac{\alpha+\gamma-1}{\alpha}
=\frac{R(\alpha+\gamma-1)}
{\alpha(\gamma-1)}}.
\]
\excheck $\alpha=0$ 时沿路径内能不变、
温度不变，此题以非零 $\Delta U$ 给出的
比式不能直接使用，需单独讨论等温极限。

\thermoenergyproblem{52}{68}
某理想气体的定容摩尔热容为 $C_V$。
在下列过程中，求过程摩尔热容随体积 $V$ 的函数：
(a) $T=T_0e^{\alpha V}$；
(b) $p=p_0e^{\alpha V}$。
\exsolution
对一摩尔理想气体，
第一定律给出
\[
C=\frac{\delta Q}{\mathrm dT}
=C_V+p\frac{\mathrm dV}{\mathrm dT}
=C_V+\frac{RT/V}{\mathrm dT/\mathrm dV}.
\]
(a) $T=T_0e^{\alpha V}$，
因此 $\mathrm dT/\mathrm dV=\alpha T$，
\[
\boxed{C_a(V)=C_V+\frac R{\alpha V}}.
\]
(b) $p=p_0e^{\alpha V}$，
状态方程 $T=pV/R$，
故
$\mathrm dT/\mathrm dV=(p/R)(1+\alpha V)$。
代入：
\[
\boxed{C_b(V)=C_V+\frac R{1+\alpha V}}.
\]
\excheck 在 $T'(V)=0$ 的特殊位置，
用 $\delta Q/\mathrm dT$ 定义的过程热容会发散；
与 $C_V$ 常数并不矛盾。

\thermoenergyproblem{53}{68}
一摩尔理想气体的绝热指数为 $\gamma$，
沿 $p=p_0+\alpha/V$ 变化，
其中 $p_0,\alpha>0$。
求：(a) 过程摩尔热容 $C(V)$；
(b) 当 $V_1\to V_2$ 时，
内能增量、气体做功和吸热量。
\exsolution
状态方程给出
$RT=pV=p_0V+\alpha$，
所以 $\mathrm dT=(p_0/R)\,\mathrm dV$，
$\mathrm dV/\mathrm dT=R/p_0$。
于是
\[
\boxed{C(V)=C_V+\frac{Rp}{p_0}
=C_p+\frac{R\alpha}{p_0V}
=\frac{\gamma R}{\gamma-1}+\frac{R\alpha}{p_0V}}.
\]
积分内能变化：
\[
\boxed{\Delta U=\frac{p_0(V_2-V_1)}{\gamma-1}}.
\]
气体做功
\[
\boxed{A=\int_{V_1}^{V_2}
\left(p_0+\frac{\alpha}{V}\right)\mathrm dV
=p_0(V_2-V_1)+\alpha\ln\frac{V_2}{V_1}}.
\]
所以
\[
\boxed{Q=\frac{\gamma p_0}{\gamma-1}(V_2-V_1)
+\alpha\ln\frac{V_2}{V_1}}.
\]
\excheck $p_0V$、$\alpha$ 均有能量量纲，
故每一项都为焦耳。

\thermoenergyproblem{54}{69}
一摩尔理想气体的定压摩尔热容为 $C_p$，
过程满足 $T=T_0+\alpha V$。
求：(a) 过程摩尔热容；
(b) 体积由 $V_1$ 增至 $V_2$ 的吸热量。
\exsolution
由状态方程
$p=RT/V=R(T_0/V+\alpha)$，
且 $\mathrm dT/\mathrm dV=\alpha$。
因此
\[
C=C_V+p\,\frac{\mathrm dV}{\mathrm dT}
=C_p+\frac{RT_0}{\alpha V},
\quad
\boxed{C(V)=C_p+\frac{RT_0}{\alpha V}}.
\]
沿过程吸热
\[
Q=\int_{V_1}^{V_2}C(V)\,\mathrm dT
=\alpha\int_{V_1}^{V_2}C(V)\,\mathrm dV.
\]
计算：
\[
\boxed{Q=C_p\alpha(V_2-V_1)+
RT_0\ln\frac{V_2}{V_1}}.
\]
\excheck 若 $\alpha=0$，
过程为等温且不能把 $C(V)$ 的表达式
当成有限量；但吸热的积分形式仍有极限
$Q=RT_0\ln(V_2/V_1)$。

\thermoenergyproblem{55}{69}
求一摩尔理想气体满足下列过程热容规律时，
$T,V$ 之间的过程方程：
(a) $C=C_V+\alpha T$；
(b) $C=C_V+\beta V$；
(c) $C=C_V+a p$。
\exsolution
对一摩尔理想气体，
$\delta Q=C_V\,\mathrm dT+p\,\mathrm dV$
及 $p=RT/V$ 给出
\[
(C-C_V)\,\mathrm dT
=\frac{RT}{V}\,\mathrm dV.
\]
(a) 代入 $C-C_V=\alpha T$，
\[
\alpha\,\mathrm dT=\frac R V\,\mathrm dV
\Longrightarrow
\boxed{V=V_0e^{\alpha(T-T_0)/R}}.
\]
(b) 若 $C-C_V=\beta V$，
\[
\beta\,\mathrm dT=\frac{RT}{V^2}\,\mathrm dV
\quad\Longrightarrow\quad
\frac{\mathrm dT}{T}=\frac R\beta\frac{\mathrm dV}{V^2}.
\]
积分得到
$\ln T=-R/(\beta V)+\text{常数}$，即
\[
\boxed{T\exp\left(\frac R{\beta V}\right)=\text{常数}}.
\]
(c) 若 $C-C_V=ap$，则
$ap\,\mathrm dT=p\,\mathrm dV$，消去 $p$：
\[
\boxed{V-aT=\text{常数}}.
\]
\excheck (b) 的关键是先除以 $T$ 再积分；
否则会错误得到 $T\propto1/V$。
所得式与原书第55(b)题答案一致。

\thermoenergyproblem{56}{69}
一摩尔理想气体的绝热指数为 $\gamma$，
某过程热容随温度满足 $C=\alpha/T$。
若气体从 $T_0$ 加热到 $\eta T_0$，
求：(a) 气体做功；(b) 以 $p,V$ 表示的过程方程。
\exsolution
由
$\delta Q=C\,\mathrm dT=(\alpha/T)\,\mathrm dT$，
$\mathrm dU=C_V\,\mathrm dT$，
第一定律给
\[
\delta A=p\,\mathrm dV
=\left(\frac{\alpha}{T}-C_V\right)\mathrm dT.
\]
因此
\[
\boxed{A=\alpha\ln\eta-C_VT_0(\eta-1)
=\alpha\ln\eta-\frac{RT_0(\eta-1)}{\gamma-1}}.
\]
另一方面 $p=RT/V$，
\[
R\,\frac{\mathrm dV}{V}
=\left(\frac{\alpha}{T^2}-\frac{C_V}{T}\right)\mathrm dT.
\]
积分得
\[
R\ln V=-\frac\alpha T-C_V\ln T+\text{常数}.
\]
所以
$V T^{1/(\gamma-1)}e^{\alpha/(RT)}=\text{常数}$。
代入 $T=pV/R$ 并对两边取
$\gamma-1$ 次方：
\[
\boxed{pV^\gamma
\exp\left(\frac{(\gamma-1)\alpha}{pV}\right)
=\text{常数}}.
\]
\excheck 若 $\alpha=0$，
过程为绝热，末式退化为
$pV^\gamma=\text{常数}$，
与热力学第一定律一致。
